\subsection{商模的希尔伯特–塞缪尔级数}

引理 5. — 设 A 为一个环，M 为 A-模，\((\mathbf{P}_{n})\)、\((\mathbf{Q}_{n})\) 是 M 上由子模组成的两个递降过滤。假设有 \(P_{n} \supset Q_{n}\) 和 \(\text{long}_{\mathbf{A}}(\mathbf{P}_{n}/\mathbf{Q}_{n}) < +\infty\) 对每个 \(n \in Z\) 成立，并且存在整数 \(n_{1}\)，使 \(Q_{n_{1}} = M\)。在 \(\mathbf{Z}((T))\) 中，有不等式

\[
\begin{array}{r l} \sum_ {n \in \mathbf {Z}} \mathrm{long} _ {\mathbf {A}} ((\mathrm{P} _ {n + 1} \cap \mathrm{Q} _ {n}) / \mathrm{Q} _ {n + 1}). & \mathrm{T} ^ {n} \leqslant \sum_ {n \in \mathbf {Z}} \mathrm{long} _ {\mathbf {A}} (\mathrm{P} _ {n + 1} / \mathrm{Q} _ {n + 1}). \\ & \mathrm{T} ^ {n} \leqslant \\ & \leqslant (1 - \mathrm{T}) ^ {- 1} \sum_ {n \in \mathbf {Z}} \mathrm{long} _ {\mathbf {A}} ((\mathrm{P} _ {n + 1} \cap \mathrm{Q} _ {n}) / \mathrm{Q} _ {n + 1}). \end{array}
\]

我们需要证明以下不等式：

\begin{enumerate}
\def\labelenumi{(\arabic{enumi})}
\setcounter{enumi}{14}
\tightlist
\item
\end{enumerate}

\[
\operatorname{long} \left(\left(\mathrm{P} _ {n + 1} \cap \mathrm{Q} _ {n}\right) / \mathrm{Q} _ {n + 1}\right) \leqslant \operatorname{long} \left(\mathrm{P} _ {n + 1} / \mathrm{Q} _ {n + 1}\right),\tag{16}
\]

\[
\operatorname{long} \left(\mathrm{P} _ {n + 1} / \mathrm{Q} _ {n + 1}\right) \leqslant \sum_ {i \leqslant n} \operatorname{long} \left(\left(\mathrm{P} _ {i + 1} \cap \mathrm{Q} _ {i}\right) / \mathrm{Q} _ {i + 1}\right).
\]

第一个不等式是显然的。另一方面，有 \(P_{n+1} \cap Q_i = P_{n+1}\) 对 \(i \leqslant n_1\) 成立，且 \(P_{n+1} \cap Q_{n+1} = Q_{n+1}\)；由此得到不等式

\[
\operatorname{long} \left(\mathrm{P} _ {n + 1} / \mathrm{Q} _ {n + 1}\right) \leqslant \sum_ {i \leqslant n} \operatorname{long} \left(\left(\mathrm{P} _ {n + 1} \cap \mathrm{Q} _ {i}\right) / \left(\mathrm{P} _ {n + 1} \cap \mathrm{Q} _ {i + 1}\right)\right).
\]

但是 A-模 \((\mathrm{P}_{n+1} \cap \mathrm{Q}_i)/(\mathrm{P}_{n+1} \cap \mathrm{Q}_{i+1})\) 同构于 \((\mathrm{P}_{i+1} \cap \mathrm{Q}_i)/\mathrm{Q}_{i+1}\) 的一个子模，因而不等式 (16) 成立。

引理 6. --- 设 A 为一个环，M 为 A-模，且 \((F_{n})\) 是 M 上由子模组成的递降过滤；假设存在整数 \(n_{1}\)，使 \(F_{n_{1}} = M\)。设 f 是 M 的自同态，\(M'\) 是其核，\(M''\) 是其余核。赋予 \(M'\) 过滤 \((F_{n}')\)，它由 \((F_{n})\) 诱导；赋予 \(M''\) 商过滤 \((F_{n}'')\)，它来自 \((F_{n})\)。假设 \(F_{n}/F_{n+1}\) 对每个 \(n \in \mathbf{Z}\) 都具有有限长度，并且存在整数 \(\delta\)，使 \(f(F_{n}) \subset F_{n+\delta}\)。令 \(\varphi\) 为由 f 诱导的 \(\delta\) 次分次自同态，其作用于分次模 \(\operatorname{gr}(M) = \bigoplus_{n} F_{n}/F_{n+1}\)。在 \(\mathbf{Z}((T))\) 的下列元素之间

\[
\mathrm{H} _ {\mathbf {M}} = \sum_ {n \in \mathbf {Z}} \operatorname{long} _ {\mathbf {A}} (\mathrm{F} _ {n} / \mathrm{F} _ {n + 1}). \mathrm{T} ^ {n}
\]

\[
\mathrm{H} _ {\mathbf {M} ^ {\prime}} = \sum_ {n \in \mathbf {Z}} \operatorname{long} _ {\mathbf {A}} (\mathrm{F} _ {n} ^ {\prime} / \mathrm{F} _ {n + 1} ^ {\prime}). \mathrm{T} ^ {n}
\]

\[
\mathrm{H} _ {\mathrm{M} ^ {\prime \prime}} = \sum_ {n \in \mathbf {Z}} \operatorname{long} _ {\mathrm{A}} (\mathrm{F} _ {n} ^ {\prime \prime} / \mathrm{F} _ {n + 1} ^ {\prime \prime}). \mathrm{T} ^ {n}
\]

\[
\mathrm{P} _ {\operatorname{Ker} (\varphi)} = \sum_ {n \in \mathbf {Z}} \operatorname{long} _ {\mathrm{A}} (\operatorname{Ker} (\varphi_ {n})). \mathrm{T} ^ {n},
\]

有不等式

\begin{enumerate}
\def\labelenumi{(\arabic{enumi})}
\setcounter{enumi}{16}
\tightlist
\item
\end{enumerate}

\[
\mathrm{H} _ {\mathbf {M} ^ {\prime}} \leqslant \mathrm{P} _ {\mathbf {K e r} (\varphi)}\tag{18}
\]

\[
(1 - \mathrm{T} ^ {\delta}). \mathrm{H} _ {\mathbf {M}} ^ {(1)} + \mathrm{T} ^ {\delta}. \mathrm{P} _ {\operatorname{Ker} (\varphi)} \leqslant \mathrm{H} _ {\mathbf {M} ^ {\prime \prime}} ^ {(1)} \leqslant (1 - \mathrm{T} ^ {\delta}). \mathrm{H} _ {\mathbf {M}} ^ {(1)} + \mathrm{T} ^ {\delta}. \mathrm{P} _ {\operatorname{Ker} (\varphi)} ^ {(1)}.
\]

M 上的子模序列 \(\mathbf{G}_{n}=f^{-1}(\mathbf{F}_{n+\delta})\) 是递降过滤，并且对于每个整数 n，有 \(F_{n}\subset G_{n}\)。

根据定义，有 \(\operatorname{Ker}(\varphi_n) = (\mathbf{G}_{n+1} \cap \mathbf{F}_n)/\mathbf{F}_{n+1}\)，因而

\[
\mathrm{P} _ {\operatorname{Ker} (\varphi)} = \sum_ {n \in \mathbf {Z}} \operatorname{long} _ {\mathrm{A}} ((\mathrm{G} _ {n + 1} \cap \mathrm{F} _ {n}) / \mathrm{F} _ {n + 1}). \mathrm{T} ^ {n}.\tag{19}
\]

对于每个 \(n\)，A-模 \((\mathbf{M}' \cap \mathbf{F}_n) / (\mathbf{M}' \cap \mathbf{F}_{n+1})\) 可识别为 \((\mathbf{G}_{n+1} \cap \mathbf{F}_n) / \mathbf{F}_{n+1}\) 的一个子模，不等式 (17) 立即由 (19) 得到。此外，根据引理 5，有

\[
\mathrm{P} _ {\operatorname{Ker} (\varphi)} \leqslant \sum_ {n \in \mathbf {Z}} \operatorname{long} _ {\mathrm{A}} \left(\mathrm{G} _ {n + 1} / \mathrm{F} _ {n + 1}\right). \mathrm{T} ^ {n} \leqslant \mathrm{P} _ {\operatorname{Ker} (\varphi)} ^ {(1)}.\tag{20}
\]

对于每个 \(n \in \mathbf{Z}\)，有 A-模的正合序列

\[
0 \longrightarrow \mathrm{G} _ {n + 1} / \mathrm{F} _ {n + 1} \longrightarrow \mathrm{M} / \mathrm{F} _ {n + 1} \xrightarrow {f _ {n}} \mathrm{M} / \mathrm{F} _ {n + \delta + 1} \longrightarrow \mathrm{M} ^ {\prime \prime} / \mathrm{F} _ {n + \delta + 1} ^ {\prime \prime} \longrightarrow 0,
\]

其中 \(f_{n}\) 是由 \(f\) 经过取商得到的。因此有

\[
\operatorname{long} _ {\mathrm{A}} \left(\mathrm{M} ^ {\prime \prime} / \mathrm{F} _ {n + \delta + 1} ^ {\prime \prime}\right) = \operatorname{long} _ {\mathrm{A}} \left(\mathrm{M} / \mathrm{F} _ {n + \delta + 1}\right) - \operatorname{long} _ {\mathrm{A}} \left(\mathrm{M} / \mathrm{F} _ {n + 1}\right) + \operatorname{long} _ {\mathrm{A}} \left(\mathrm{G} _ {n + 1} / \mathrm{F} _ {n + 1}\right).
\]

两边乘以 \(T^{n+\delta}\) 并对 n 求和，得到

\[
\mathrm{H} _ {\mathbf {M} ^ {\prime \prime}} ^ {(1)} = (1 - \mathrm{T} ^ {\delta}) \mathrm{H} _ {\mathbf {M}} ^ {(1)} + \mathrm{T} ^ {\delta}. \sum_ {n \in \mathbf {Z}} \operatorname{long} _ {\mathrm{A}} (\mathrm{G} _ {n + 1} / \mathrm{F} _ {n + 1}). \mathrm{T} ^ {n},\tag{21}
\]

不等式 (18) 立即由 (20) 和 (21) 得到。

引理 7. — 保持引理 6 的记号。

\begin{enumerate}
\def\labelenumi{\alph{enumi})}
\item
  有不等式 \(\mathbf{H}_{\mathbf{M}^{\prime \prime}}^{(1)}\geqslant \frac{1 - \mathrm{T}^{\delta}}{1 - \mathrm{T}}\mathbf{H}_{\mathbf{M}}.\)
\item
  等号成立的充要条件是 φ 为单射。
\item
  在此情况下，有 \(\mathbf{M}^{\prime}\subset \bigcap_{n}\mathbf{F}_{n}\)，且 A-模序列
\end{enumerate}

\[
0 \longrightarrow \operatorname{gr} (\mathrm{M}) \stackrel {\varphi} {\longrightarrow} \operatorname{gr} (\mathrm{M}) \stackrel {v} {\longrightarrow} \operatorname{gr} (\mathrm{M} ^ {\prime \prime}) \longrightarrow 0,
\]

其中 v 是典范映射，是正合的。

断言 a) 和 b) 由引理 6 的公式 (18) 以及定义 \(\mathbf{H}_{\mathbf{M}}^{(1)} = (1 - \mathbf{T})^{-1}.\mathbf{H}_{\mathbf{M}}\) 得到。

假设 \(\varphi\) 是单射。由 III, § 2, no. 8, Th. 1, (i)，有

\[
\operatorname{Ker} (f) \subset f ^ {- 1} (\mathrm{F} _ {n + \delta}) = \mathrm{F} _ {n}
\]

对于所有 n，从而得到 c) 的第一个断言。此外，有正合序列

\[
0 \longrightarrow \mathrm{M} / \mathrm{M} ^ {\prime} \stackrel {f ^ {\prime}} {\longrightarrow} \mathrm{M} \longrightarrow \mathrm{M} / f (\mathrm{M}) \longrightarrow 0,
\]

其中 \(f'\) 是由 \(f\) 经过取商得到的。若 \(\varphi\) 是单射，如上有 \(f'^{-1}(\mathbf{F}_n) = \mathbf{F}_{n-\delta}/\mathbf{M}'\)。在 \(\mathbf{M}/\mathbf{M}'\) 上，由 \(f'\) 的逆像从过滤 \(\mathbf{F}\)（定义在 \(\mathbf{M}\) 上）导出的过滤是 \(n \mapsto \mathbf{F}_{n-\delta}/\mathbf{M}'\)；其伴随分次模为 \(\operatorname{gr}(\mathbf{M})(-\delta)\)，并且有分次模的正合序列（III, § 2, no. 4, prop. 2）

\[
0 \longrightarrow \operatorname{gr} (\mathrm{M}) (- \delta) \xrightarrow {\varphi^ {\prime}} \operatorname{gr} (\mathrm{M}) \longrightarrow \operatorname{gr} \left(\mathrm{M} ^ {\prime \prime}\right) \longrightarrow 0,
\]

其中 \(\varphi_{n}^{\prime} = \varphi_{n - \delta}\) 对每个 \(n\) 都成立。这完成了 \(c\) ) 的证明。

命题 6. — 设 A 为诺特环，M 为有限生成 A-模，q 为 A 的理想，且 M/qM 具有有限长度。设 F 是 M 的 q-良过滤，令 gr(A) = \(\bigoplus_{n\geqslant 0}(\mathfrak{q}^n /\mathfrak{q}^{n + 1})\) 为 A 关于 q-进过滤的伴随分次环。

令 \((x_{1},\ldots,x_{s})\) 为 A 中元素组成的序列，\((\delta_{1},\ldots,\delta_{s})\) 为严格正整数序列，使得 \(x_{i}\in\mathfrak{q}^{\delta_{i}}\) 对 \(1\leqslant i\leqslant s\) 成立，并令 \(\xi_{i}\) 为 \(x_{i}\) 在 \(\mathrm{gr}_{\delta_{i}}(\mathbf{A})=\mathfrak{q}^{\delta_{i}}/\mathfrak{q}^{\delta_{i}+1}\) 中的类。

\begin{enumerate}
\def\labelenumi{\alph{enumi})}
\tightlist
\item
  赋予 A-模 \(\overline{\mathbf{M}} = \mathbf{M}/(x_1\mathbf{M} + \cdots + x_s\mathbf{M})\) F 的商 q-良过滤 \(\overline{\mathbf{F}}\)。于是，在 \(\mathbf{Z}((\mathbf{T}))\) 中有不等式
\end{enumerate}

\[
\mathrm{H} _ {\mathrm{M}, \overline {{\mathrm{F}}}} ^ {(s)} \geqslant \left(\prod_ {i = 1} ^ {s} \frac {1 - \mathrm{T} ^ {\delta_ {i}}}{1 - \mathrm{T}}\right). \mathrm{H} _ {\mathrm{M}, \mathrm{F}}.\tag{22}
\]

\begin{enumerate}
\def\labelenumi{\alph{enumi})}
\setcounter{enumi}{1}
\tightlist
\item
  (22) 中取等号的充要条件是：环 gr(A) 的元素序列 \((\xi_1, \ldots, \xi_s)\) 是模 gr(M) = \(\bigoplus_{n} (\mathbf{F}_n / \mathbf{F}_{n+1})\) 的完全截切序列。
\end{enumerate}

在此情形，从 \(\mathrm{gr}(\mathbf{M}) / \sum_{i=1}^{s} \xi_i \cdot \mathrm{gr}(\mathbf{M})\) 到 \(\mathrm{gr}(\overline{\mathbf{M}}) = \bigoplus_{n} (\overline{\mathbf{F}}_n / \overline{\mathbf{F}}_{n+1})\) 的典范同态是同构。

\begin{enumerate}
\def\labelenumi{\alph{enumi})}
\setcounter{enumi}{2}
\tightlist
\item
  假设 b) 的条件满足，并且每个 A-模 \(\mathbf{M}_{i} = \mathbf{M}/(x_{1}\mathbf{M} + \cdots + x_{i}\mathbf{M})\) 对 \(0 \leqslant i < s\) 都关于 q-进拓扑分离 \footnote{当 \(q\) 包含于 \(A\) 的根基中时尤其如此（III, § 3, no. 3, Prop. 6）。}。则序列 \((x_{1}, ..., x_{s})\) 是 A-模 M 的完全截切序列。
\end{enumerate}

当 \(s = 1\) 时，有 \(\bigcap_{n} F_{n} = \bigcap_{n} q^{n} M\)，且序列 \(\{\xi_{1}\}\) 是 \(\operatorname{gr}(M)\) 的完全截切序列，当且仅当比率为 \(\xi_{1}\) 的倍乘映射在 \(\operatorname{gr}(M)\) 上是单射。将引理 7 应用于倍乘映射 \(f = (x_{1})_{M}\)，它作用于 \(M\)，即可立即得到命题 6。

现在设 \(s \geqslant 2\)，并对 \(s\) 归纳论证。将归纳假设应用于赋予 F 的商过滤 G 的 A-模 \(\mathbf{M}_{1} = \mathbf{M}/x_{1}\mathbf{M}\)，以及序列 \((x_{2}, ..., x_{s})\)，得到不等式

\[
\mathrm{H} _ {\overline {{\mathrm{M,F}}}} ^ {(s - 1)} \geqslant \left(\prod_ {i = 2} ^ {s} \frac {1 - \mathrm{T} ^ {\delta_ {i}}}{1 - \mathrm{T}}\right). \mathrm{H} _ {\mathrm{M} _ {1}, \mathrm{G}};\tag{23}
\]

等号成立当且仅当序列 \((\xi_{2}, \ldots, \xi_{s})\) 是 \(\operatorname{gr}(A)\)-模 \(\operatorname{gr}(\mathbf{M}_{1}) = \bigoplus_{n} \mathrm{G}_{n}/\mathrm{G}_{n+1}\) 的完全截切序列。由于元素 \(\frac{1 - \mathrm{T}^{\delta_{i}}}{1 - \mathrm{T}}\) 属于 \(\mathbf{Z}((\mathrm{T}))\) 且为正，已处理的 \(s = 1\) 情形和公式 (23) 给出不等式

\[
\mathrm{H} _ {\mathrm{M}, \mathrm{F}} ^ {(s)} \geqslant \left(\prod_ {i = 2} ^ {s} \frac {1 - \mathrm{T} ^ {\delta_ {i}}}{1 - \mathrm{T}}\right). \mathrm{H} _ {\mathrm{M} _ {1}, \mathrm{G}} ^ {(1)} \geqslant \left(\prod_ {i = 1} ^ {s} \frac {1 - \mathrm{T} ^ {\delta_ {i}}}{1 - \mathrm{T}}\right). \mathrm{H} _ {\mathrm{M}, \mathrm{F}}.\tag{24}
\]

这证明了 a)。

只有在 (23) 中同时取等号且有等式

\[
\mathrm{H} _ {\mathbf {M} _ {1}, \mathbf {G}} ^ {(1)} = \left(\frac {1 - \mathrm{T} ^ {\delta_ {1}}}{1 - \mathrm{T}}\right) \cdot \mathrm{H} _ {\mathbf {M}, \mathbf {F}}.\tag{25}
\]

后一关系意味着 \(\{\xi_{1}\}\) 是 \(\operatorname{gr}(M)\) 的完全截切序列，并蕴含从 \(\operatorname{gr}(M)/\xi_{1}\operatorname{gr}(M)\) 到 \(\operatorname{gr}(M_{1})\) 的典范同态是同构。换言之，(22) 中取等号，当且仅当 \(\{\xi_{1}\}\) 是 \(\operatorname{gr}(M)\) 的完全截切序列，且 \(\{\xi_{2}, ..., \xi_{s}\}\) 是 \(\operatorname{gr}(M)/\xi_{1}\operatorname{gr}(M)\) 的完全截切序列。这意味着 \(\{\xi_{1}, ..., \xi_{s}\}\) 是 \(\operatorname{gr}(M)\) 的完全截切序列（A, X, p. 160, cor. 2）。因此证明了 \(b\)) 的两个条件等价。假设它们满足；则根据归纳假设，\(\operatorname{gr}(M)\) 可识别为 \(\operatorname{gr}(M_{1})/\sum_{i=2}^{s} \xi_{i}\operatorname{gr}(M_{1})\)；此外，由于 \(\operatorname{gr}(M_{1})\) 可识别为 \(\operatorname{gr}(M)/\xi_{1}\operatorname{gr}(M)\)，\(b\)) 的最后一个断言也得到满足。

现在假设 \(\{\xi_1, \ldots, \xi_s\}\) 是 \(\operatorname{gr}(\mathbf{M})\) 的完全截切序列，且每个 \(\mathbf{M}_i\) 关于 q-进拓扑都是分离的（\(0 \leqslant i < s\)）。由上面的结果和归纳假设，序列 \((x_2, \ldots, x_s)\) 是 \(\mathbf{M}_1\) 的完全截切序列；由于有 \(\mathbf{M}_1 = \mathbf{M}/x_1\mathbf{M}\)，且 \(\{x_1\}\) 是 \(\mathbf{M}\) 的完全截切序列，序列 \((x_1, x_2, \ldots, x_s)\) 是 \(\mathbf{M}\) 的完全截切序列（A, X, p. 160, th. 1）。

